% Unit 052 terjemahan Bahasa Indonesia dari chapter4.tex baris 561--622.
% Sumber: Methods of Algebra, Volume 2, commit
% 9a5803ff77dd3257484cb177f851a73770a59dd3 (CC BY 4.0).
% stable-unit-id: o014.aljabr2.chapter4.triangulated-localization
% Cakupan: awal Bagian 4.3, dari motivasi dan sistem multiplikatif yang
% kompatibel dengan triangulasi sampai lokalisasi yang mempertahankan struktur
% (pra)bertriangulasi; berhenti sebelum konstruksi $S\mathcal{N}$ pada baris 623.

% segment-id: o014.li2.chapter4.triangulated-localization.h001
\section{Lokalisasi Kategori Bertriangulasi}\label{sec:triangulated-localization}
\hypertarget{o014.aljabr2.chapter4.triangulated-localization}{ }

% segment-id: o014.li2.chapter4.triangulated-localization.p001
Bagian ini mengkaji hubungan antara kategori bertriangulasi dan teori
lokalisasi Gabriel--Zisman. Secara khusus,
\begin{compactitem}
	\item konstruksi lokalisasi kategori dalam \S\ref{sec:cat-localization}
	secara formal menambahkan invers bagi morfisme-morfisme dalam suatu sistem
	multiplikatif, dan tugas pertama bagian ini adalah menunjukkan bahwa
	konstruksi tersebut mempertahankan struktur (pra)bertriangulasi;
	\item dalam kategori bertriangulasi, lokalisasi juga dapat dipandang sebagai
	proses menolkan suatu subkategori bertriangulasi secara formal; gagasan ini
	serupa dengan konstruksi kuosien Serre dalam \S\ref{sec:Serre-subcat}.
\end{compactitem}

% segment-id: o014.li2.chapter4.triangulated-localization.p002
Selanjutnya, $(\mathcal{D}, T)$ selalu dianggap sebagai kategori
prabertriangulasi. Jika pembahasan diperluas dari kategori aditif ke konteks
$\Bbbk$-linear,
	dengan $\Bbbk$ suatu gelanggang komutatif sembarang, semua hasil dalam bagian
	ini mempunyai generalisasi yang sepadan; perinciannya tidak diberikan.

% segment-id: o014.li2.chapter4.triangulated-localization.d001
\begin{definisi}
	\index{sistem multiplikatif!kompatibel dengan triangulasi}
	Misalkan $S \subset \Mor(\mathcal{D})$ suatu sistem multiplikatif dari
	Definisi~\ref{def:multiplicative-family}. Sistem $S$ disebut kompatibel
	dengan triangulasi jika kondisi-kondisi berikut dipenuhi.
	\begin{enumerate}[(ST1)]
		\item Untuk $s \in \Mor(\mathcal{D})$, berlaku $s \in S$ jika dan hanya
		jika $Ts \in S$.
		\item Tinjau diagram berikut, yang bagian bergaris penuhnya diberikan:
% segment-id: o014.li2.chapter4.triangulated-localization.g001
		\[\begin{tikzcd}
			X \arrow[d, "\alpha"'] \arrow[r] & Y \arrow[d, "\beta"] \arrow[r] & Z \arrow[dashed, d, "\gamma"] \arrow[r] & TX \arrow[dashed, d, "T\alpha"] \\
			X' \arrow[r] & Y' \arrow[r] & Z' \arrow[r] & TX'
		\end{tikzcd}\]
		Jika $\alpha, \beta \in S$, maka terdapat morfisme $\gamma$ yang ditunjukkan
		dengan garis putus-putus sedemikian sehingga seluruh diagram komutatif dan
		$\gamma \in S$.
	\end{enumerate}
\end{definisi}

% segment-id: o014.li2.chapter4.triangulated-localization.p003
Perhatikan bahwa (ST2) merupakan penguatan (TR4) bagi $S$.

% segment-id: o014.li2.chapter4.triangulated-localization.pr001
\begin{proposisi}\label{prop:localization-triangulated}
	Misalkan suatu sistem multiplikatif $S \subset \Mor(\mathcal{D})$ kompatibel
	dengan triangulasi. Maka lokalisasi $\mathcal{D}[S^{-1}]$ mempunyai struktur
	kategori prabertriangulasi yang unik sedemikian sehingga
	$Q: \mathcal{D} \to \mathcal{D}[S^{-1}]$ merupakan funktor bertriangulasi.
	Hingga isomorfisme, segitiga-segitiga terbedakan dalam
	$\mathcal{D}[S^{-1}]$ tepat merupakan citra segitiga-segitiga terbedakan
	dalam $\mathcal{D}$. Jika $\mathcal{D}$ kategori bertriangulasi, maka
	$\mathcal{D}[S^{-1}]$ juga demikian.
\end{proposisi}

% segment-id: o014.li2.chapter4.triangulated-localization.pf001
\begin{bukti}
	Teorema~\ref{prop:localization-additivity} menunjukkan bahwa
	$\mathcal{D}[S^{-1}]$ merupakan kategori aditif. Selain itu, (ST1) dan sifat
	universal lokalisasi memberikan sepasang funktor unik dari
	$\mathcal{D}[S^{-1}]$ ke dirinya sendiri yang membuat diagram berikut
	komutatif:
% segment-id: o014.li2.chapter4.triangulated-localization.g002
	\[\begin{tikzcd}
		\mathcal{D} \arrow[r, "T", shift left] \arrow[d, "Q"'] & \mathcal{D} \arrow[l, "{T^{-1}}", shift left] \arrow[d, "Q"] \\
		\mathcal{D}[S^{-1}] \arrow[r, shift left] & \mathcal{D}[S^{-1}] \arrow[l, shift left]
	\end{tikzcd}\]
	Sifat universal juga menunjukkan bahwa kedua funktor pada baris kedua saling
	invers. Untuk menyederhanakan notasi, keduanya tetap ditulis sebagai
% segment-id: o014.li2.chapter4.triangulated-localization.g003
	$\begin{tikzcd}
		\mathcal{D}[S^{-1}] \arrow[r, "T", shift left] & \mathcal{D}[S^{-1}] \arrow[l, "{T^{-1}}", shift left]
	\end{tikzcd}$.
	Dengan demikian, $(\mathcal{D}[S^{-1}], T)$ menjadi kategori aditif dengan
	pergeseran.
	
	Kita terlebih dahulu membuktikan keunikan struktur prabertriangulasi.
	Menurut definisi, funktor bertriangulasi $Q$ harus mempertahankan
	segitiga-segitiga terbedakan. Sebaliknya, misalkan
	$X \xrightarrow{f} Y \to Z \xrightarrow{+1}$ suatu segitiga terbedakan dalam
	$\mathcal{D}[S^{-1}]$. Setelah menggantinya dengan segitiga yang isomorfik, kita boleh
	mengandaikan bahwa $f: X \to Y$ berasal dari suatu morfisme dalam
	$\mathcal{D}$, yang demi menghindari kerancuan ditulis
	$f_0: X_0 \to Y_0$. Dengan (TR2), perluas $f_0$ menjadi segitiga terbedakan
	$X_0 \xrightarrow{f_0} Y_0 \to Z_0 \xrightarrow{+1}$. Maka
	Korolari~\ref{prop:Cone-uniqueness} menjamin bahwa citranya di bawah $Q$
	isomorfik dengan segitiga semula.
	
	Sekarang kita membuktikan keberadaan struktur prabertriangulasi. Definisikan
	segitiga-segitiga terbedakan dalam $\mathcal{D}[S^{-1}]$ seperti yang
	dinyatakan dalam proposisi. Aksioma (TR0) langsung terpenuhi, sedangkan (TR1)
	dan (TR3) diwarisi langsung dari $\mathcal{D}$. Untuk (TR2), tinjau suatu
	morfisme $f: X \to Y$ dalam $\mathcal{D}[S^{-1}]$. Seperti biasa, setelah
	mengganti objek-objeknya melalui isomorfisme yang berasal dari $S$, kita boleh
	mengandaikan bahwa $f$ berasal dari $\mathcal{D}$; dengan demikian (TR2) juga
	cukup menggunakan pernyataan dalam $\mathcal{D}$.
	
	Untuk (TR4), tinjau diagram komutatif berikut dalam
	$\mathcal{D}[S^{-1}]$; bagian bergaris putus-putus akan dibahas kemudian.
% segment-id: o014.li2.chapter4.triangulated-localization.g004
	\begin{equation}\label{eqn:localization-triangulated-aux}\begin{tikzcd}
		X \arrow[r, "f"] \arrow[d, "\alpha"'] & Y \arrow[r, "g"] \arrow[d, "\beta"] & Z \arrow[r, "h"] \arrow[dashed, d, "\gamma"] & TX \arrow[dashed, d, "T\alpha"] \\
		X' \arrow[r, "{f'}"'] & Y' \arrow[r, "{g'}"'] & Z' \arrow[r, "{h'}"'] & T X'
	\end{tikzcd}\end{equation}
	Di sini setiap baris merupakan segitiga terbedakan dan semua panah horizontal
	berasal dari $\mathcal{D}$. Berdasarkan konstruksi dalam
	\S\ref{sec:cat-localization}, dapat dibuktikan bahwa terdapat
	$A, B \in \Obj(\mathcal{D})$, suatu morfisme
	$[A \to B] \in \Mor(\mathcal{D})$, serta $s, t \in S$, sedemikian sehingga
	$\alpha = as^{-1}$ dan $\beta = bt^{-1}$ dalam $\mathcal{D}[S^{-1}]$ dan
	bagian bergaris penuh pada diagram berikut komutatif dalam $\mathcal{D}$;
	bagian bergaris putus-putus dan objek $C$ akan dibahas kemudian.
% segment-id: o014.li2.chapter4.triangulated-localization.g005
	\[\begin{tikzcd}
		X \arrow[r, "f"] & Y \arrow[r, "g"] & Z \arrow[r, "h"] & TX \\
		A \arrow[u, "s"] \arrow[d, "a"'] \arrow[r] & B \arrow[u, "t"] \arrow[d, "b"'] \arrow[dashed, r] & C \arrow[dashed, u, "u"] \arrow[dashed, r] \arrow[dashed, d, "c"'] & TA \arrow[dashed, u, "Ts"'] \arrow[dashed, d, "Ta"] \\
		X' \arrow[r, "{f'}"'] & Y' \arrow[r, "{g'}"'] & Z' \arrow[r, "{h'}"'] & T X' .
	\end{tikzcd}\]
	Konstruksi eksplisitnya agak rumit dan dijadikan latihan berpanduan dalam
	\CHref{sec:categories}.

	Sekarang terapkan (TR2) pada $A \to B$ dalam $\mathcal{D}$ untuk memperoleh
	segitiga terbedakan $A \to B \to C \xrightarrow{+1}$ pada baris kedua.
	Kemudian terapkan (ST2) untuk memperoleh $C \xrightarrow{u \in S} Z$, lalu
	terapkan (TR4) untuk memperoleh $C \xrightarrow{c} Z'$, sedemikian sehingga
	seluruh diagram komutatif dalam $\mathcal{D}$. Dalam
	$\mathcal{D}[S^{-1}]$, ambil $\gamma := c u^{-1}$. Dengan pilihan ini,
	seluruh diagram \eqref{eqn:localization-triangulated-aux} komutatif. Jadi (TR4)
	telah diverifikasi.
	
	Jika $\mathcal{D}$ merupakan kategori bertriangulasi, (TR5) untuk
	$\mathcal{D}[S^{-1}]$ juga cukup diverifikasi dalam $\mathcal{D}$.
\end{bukti}
